\section{Statement of the improvement}

Ge--Lin prove hyperbolic circle packing existence from the pointwise
condition
\[
  L_i(\T,\Phiw)>2\pi \qquad (i\in V).
\]
The pointwise condition can be replaced by a corrected average condition
over vertex subsets.  This keeps Chow--Luo's hyperbolic combinatorial Ricci
flow as the main tool.  The correction term below is exactly the boundary
term appearing in Chow--Luo's curvature image inequalities, rewritten in
the language of Ge--Lin's character.

\section{Characters of subsets}

Let $(X,\T,\Phiw)$ be a weighted triangulated closed surface with
$\Phiw:E\to[0,\pi/2]$.  For a face $f=ijk$, define
\[
  \gamma_i^f
  =
  \arccos
  \frac{1+\cos\Phiw_{ij}+\cos\Phiw_{ik}-\cos\Phiw_{jk}}
       {2\sqrt{1+\cos\Phiw_{ij}}\sqrt{1+\cos\Phiw_{ik}}}.
\]
Thus
\[
  L_i(\T,\Phiw)=\sum_{f\ni i}\gamma_i^f .
\]
These numbers are not the angles of the hyperbolic triangles appearing
along the Ricci flow.  They are the limiting comparison angles in
Ge--Lin's definition of the character, equivalently the angles of the
Euclidean equal-radius comparison triangle.  We record this as a lemma to
avoid ambiguity.

\begin{lemma}[comparison-angle identity]\label{lem:comparison-angle}
\lean{ProbabilityCharacter.comparison_angle_sum_eq_pi}\leanok
For every face $f=ijk$,
\[
  \gamma_i^f+\gamma_j^f+\gamma_k^f=\pi .
\]
\end{lemma}

\begin{proof}
Consider the Euclidean comparison triangle with edge lengths
\[
  \ell_{ij}=\sqrt{2+2\cos\Phiw_{ij}},\quad
  \ell_{ik}=\sqrt{2+2\cos\Phiw_{ik}},\quad
  \ell_{jk}=\sqrt{2+2\cos\Phiw_{jk}}.
\]
These are the Euclidean circle-packing edge lengths with all three radii
equal to one.  The triangle inequalities for these lengths are exactly the
Euclidean-background version of Thurston's triangle-inequality lemma used
by Ge--Lin.

Let $\alpha_i$ be the Euclidean angle opposite the edge $\ell_{jk}$.
By the Euclidean cosine law,
\[
 \cos \alpha_i
 =
 \frac{\ell_{ij}^2+\ell_{ik}^2-\ell_{jk}^2}{2\ell_{ij}\ell_{ik}}
 =
 \frac{1+\cos\Phiw_{ij}+\cos\Phiw_{ik}-\cos\Phiw_{jk}}
      {2\sqrt{1+\cos\Phiw_{ij}}\sqrt{1+\cos\Phiw_{ik}}}.
\]
Since $0<\alpha_i<\pi$, this gives $\alpha_i=\gamma_i^f$.  Similarly
$\alpha_j=\gamma_j^f$ and $\alpha_k=\gamma_k^f$.  Hence the sum is the
Euclidean triangle angle sum $\pi$.
\end{proof}

For $I\subseteq V$, let $\mathcal F_m(I)$ denote the set of faces having
exactly $m$ vertices in $I$.  Put
\[
  L_\Phi(I)=\sum_{i\in I}L_i .
\]
If $f\in\mathcal F_2(I)$, let $\gamma_{\mathrm{out}}^f$ be the comparison
angle at the unique vertex of $f$ outside $I$.  If $f\in\mathcal F_1(I)$,
let $v_f$ be the unique vertex of $f$ in $I$, and let $e_f$ be the opposite
edge of $f$, whose two endpoints are outside $I$.

\begin{definition}[finite combinatorial model]\label{def:finite-combinatorial-model}
\lean{
ProbabilityCharacter.Combinatorics.TriFace,
ProbabilityCharacter.Combinatorics.TriFace.internalVertexPairs,
ProbabilityCharacter.Combinatorics.TriFace.card_internalVertexPairs,
ProbabilityCharacter.Combinatorics.TriFace.choose_two_countIn_eq_piecewise,
ProbabilityCharacter.Combinatorics.TriEdge,
ProbabilityCharacter.Combinatorics.TriEdge.ext,
ProbabilityCharacter.Combinatorics.TriEdge.vertsEmbedding,
ProbabilityCharacter.Combinatorics.TriangulationModel,
ProbabilityCharacter.Combinatorics.TriangulationModel.F1,
ProbabilityCharacter.Combinatorics.TriangulationModel.F2,
ProbabilityCharacter.Combinatorics.TriangulationModel.F3,
ProbabilityCharacter.Combinatorics.TriangulationModel.internalEdges,
ProbabilityCharacter.Combinatorics.TriangulationModel.internalEdgesInFace,
ProbabilityCharacter.Combinatorics.TriangulationModel.internalEdgeVertexSetsInFace,
ProbabilityCharacter.Combinatorics.TriangulationModel.card_internalEdgeVertexSetsInFace,
ProbabilityCharacter.Combinatorics.TriangulationModel.internalEdgesInFace_card_eq_choose_of_vertexSet_image,
ProbabilityCharacter.Combinatorics.TriangulationModel.edgeFaceIncidences,
ProbabilityCharacter.Combinatorics.TriangulationModel.edgeSideIncidenceCount,
ProbabilityCharacter.Combinatorics.TriangulationModel.faceSideIncidenceCount,
ProbabilityCharacter.Combinatorics.TriangulationModel.edgeFaceIncidences_card_eq_edgeSideIncidenceCount,
ProbabilityCharacter.Combinatorics.TriangulationModel.edgeFaceIncidences_card_eq_faceSideIncidenceCount,
ProbabilityCharacter.Combinatorics.TriangulationModel.edgeSideIncidenceCount_eq_faceSideIncidenceCount,
ProbabilityCharacter.Combinatorics.TriangulationModel.edgeSideIncidenceCount_eq_two_mul,
ProbabilityCharacter.Combinatorics.TriangulationModel.faceSideIncidenceCount_eq_three_mul_F3_add_F2_of_face_contribution,
ProbabilityCharacter.Combinatorics.TriangulationModel.faceSideIncidenceCount_eq_three_mul_F3_add_F2_of_choose_contribution,
ProbabilityCharacter.Combinatorics.TriangulationModel.subsetEdgeCountIdentity_of_internalEdgeVertexSet_images,
ProbabilityCharacter.Combinatorics.TriangulationModel.subsetEdgeCountIdentity_of_choose_face_contribution,
ProbabilityCharacter.Combinatorics.TriangulationModel.subsetEdgeCountIdentity_of_local_face_contribution,
ProbabilityCharacter.Combinatorics.TriangulationModel.subsetEdgeCountIdentity_of_face_side_count,
ProbabilityCharacter.Combinatorics.TriangulationModel.subsetEdgeCountIdentity_of_incidence_counts,
ProbabilityCharacter.Combinatorics.TriangulationModel.subsetEdgeCountIdentity,
ProbabilityCharacter.Combinatorics.TriangulationModel.subsetEulerCount}\leanok
This is the finite bookkeeping model used in the obstruction identity.  It
separates the purely combinatorial part of the proof from the geometric input
about closed triangulated surfaces.

In Lean, a face is represented by a three-element finite set of vertices and
an edge by a two-element finite set.  A finite triangulation model consists of
finite sets of vertices, edges, and triangular faces.  For a vertex subset
\(I\), the model defines:
\[
  \mathcal F_1(I),\quad \mathcal F_2(I),\quad \mathcal F_3(I),\quad
  E_I,\quad |I|-|E_I|+|\mathcal F_3(I)|.
\]
Here \(\mathcal F_m(I)\) is the set of faces with exactly \(m\) vertices in
\(I\), and \(E_I\) is the set of edges whose endpoints both lie in \(I\).

The formalized counting argument has two parts.  First, Lean double-counts
edge--face incidences, giving an edge-side count and a face-side count.
Second, Lean reduces the local contribution of a single face to the number
of two-element subsets of \(f\cap I\).  This local number is \(3\), \(1\), or
\(0\), according as \(f\in\mathcal F_3(I)\), \(f\in\mathcal F_2(I)\), or
otherwise.  Thus the model proves the combinatorial identity
\[
  2|E_I|=3|\mathcal F_3(I)|+|\mathcal F_2(I)|
\]
once the usual closed-surface input is supplied: every internal edge is
incident to exactly two faces.
\end{definition}

\begin{definition}[boundary-corrected average character]\label{def:boundary-corrected-average}
\lean{ProbabilityCharacter.correctedAverage}\leanok
For $\emptyset\ne I\subsetneq V$, define
\[
  \Delta_\Phi(I)
  =
  \sum_{f\in\mathcal F_2(I)}\gamma_{\mathrm{out}}^f
  +
  \sum_{f\in\mathcal F_1(I)}
    \bigl(\pi-\Phiw(e_f)-\gamma_{v_f}^f\bigr),
\]
and
\[
  \ac_\Phi(I)=\frac{L_\Phi(I)+\Delta_\Phi(I)}{|I|}.
\]
For $I=V$, set $\Delta_\Phi(V)=0$ and
\[
  \ac_\Phi(V)=\frac{L_\Phi(V)}{|V|}.
\]
\end{definition}

\begin{lemma}[non-negativity of the correction]\label{lem:nonneg-correction}
\uses{def:boundary-corrected-average,lem:comparison-angle}
\lean{ProbabilityCharacter.DeltaPhi_nonneg_of_valid_angles}\leanok
For every nonempty $I\subseteq V$,
\[
  \Delta_\Phi(I)\ge 0.
\]
\end{lemma}

\begin{proof}
The summands $\gamma_{\mathrm{out}}^f$ are non-negative.  It remains to
check a summand of the second type.  Let $f=ijk$ and suppose $i=v_f$ while
$e_f=jk$.  Since $\Phiw:E\to[0,\pi/2]$, the formula for $\gamma_i^f$ gives
\[
 \cos \gamma_i^f
 =
 \frac{1+\cos\Phiw_{ij}+\cos\Phiw_{ik}-\cos\Phiw_{jk}}
      {2\sqrt{1+\cos\Phiw_{ij}}\sqrt{1+\cos\Phiw_{ik}}}
 \ge 0,
\]
because $1-\cos\Phiw_{jk}\ge0$ and the other cosine terms are
non-negative.  Hence $\gamma_i^f\le\pi/2$.  Also
$\pi-\Phiw_{jk}\ge\pi/2$.  Therefore
\[
  \pi-\Phiw(e_f)-\gamma_{v_f}^f
  =
  \pi-\Phiw_{jk}-\gamma_i^f
  \ge 0.
\]
\end{proof}

\section{The exact obstruction identity}

For $\emptyset\ne I\subsetneq V$, let $F_I$ be the full subcomplex of
$\T$ spanned by $I$.  Following Chow--Luo, define
\[
  B_\Phi(I)
  =
  \sum_{(e,v)\in\Lk(I)}(\pi-\Phiw(e))-2\pi\chi(F_I),
\]
where $\Lk(I)$ is the set of pairs $(e,v)$ such that the endpoints of $e$
are not in $I$, $v\in I$, and $e$ together with $v$ forms a triangle of
$\T$.

\begin{lemma}[obstruction equals corrected surplus]\label{lem:obstruction-corrected-surplus}
\uses{def:finite-combinatorial-model,def:boundary-corrected-average}
\lean{ProbabilityCharacter.obstruction_eq_card_mul_surplus, ProbabilityCharacter.Combinatorics.obstruction_identity_from_count_data}\leanok
For every nonempty proper subset $I\subsetneq V$,
\[
  B_\Phi(I)=L_\Phi(I)-2\pi |I|+\Delta_\Phi(I)
  =
  |I|\bigl(\ac_\Phi(I)-2\pi\bigr).
\]
\end{lemma}

\begin{proof}
Let $E_I$ be the set of edges whose endpoints both lie in $I$.  Let
$F_3=\mathcal F_3(I)$ and $F_2=\mathcal F_2(I)$.  Since $X$ is closed,
every edge of $E_I$ is incident to two faces.  Counting incidences between
edges of $E_I$ and adjacent faces gives
\[
  2|E_I|=3|F_3|+|F_2|.
\]
The subset character is
\[
  L_\Phi(I)
  =
  \pi|F_3|
  +\sum_{f\in\mathcal F_2(I)}(\pi-\gamma_{\mathrm{out}}^f)
  +\sum_{f\in\mathcal F_1(I)}\gamma_{v_f}^f .
\]
After adding $\Delta_\Phi(I)$, the terms on $\mathcal F_2(I)$ and
$\mathcal F_1(I)$ simplify to
\[
  L_\Phi(I)+\Delta_\Phi(I)
  =
  \pi|F_3|+\pi|F_2|
  +\sum_{f\in\mathcal F_1(I)}(\pi-\Phiw(e_f)).
\]
Using $2|E_I|=3|F_3|+|F_2|$, this becomes
\[
  L_\Phi(I)+\Delta_\Phi(I)
  =
  2\pi |E_I|-2\pi |F_3|
  +\sum_{f\in\mathcal F_1(I)}(\pi-\Phiw(e_f)).
\]
Since
\[
  \chi(F_I)=|I|-|E_I|+|F_3|,
\]
we get
\[
  L_\Phi(I)-2\pi |I|+\Delta_\Phi(I)
  =
  \sum_{f\in\mathcal F_1(I)}(\pi-\Phiw(e_f))-2\pi\chi(F_I).
\]
Each face in $\mathcal F_1(I)$ contributes exactly one pair $(e_f,v_f)$ to
$\Lk(I)$, and every pair in $\Lk(I)$ arises this way.  Hence the last
expression is $B_\Phi(I)$.
\end{proof}

\section{Average criterion}

\begin{theorem}[Chow--Luo, zero-curvature case]\label{thm:chow-luo}
For a weighted triangulated closed surface $(X,\T,\Phiw)$ in the hyperbolic
background, there is a radius vector $r_{ze}\in\R_{>0}^V$ with
$K(r_{ze})=0$ if and only if $\chi(X)<0$ and
\[
  B_\Phi(I)>0
  \qquad
  (\emptyset\ne I\subsetneq V).
\]
In this case the hyperbolic combinatorial Ricci flow
\[
  \frac{dr_i}{dt}=-K_i\sinh r_i
\]
exists for all time and converges exponentially fast, from every initial
radius vector in $\R_{>0}^V$, to $r_{ze}$.  The vector $r_{ze}$ is unique.
\end{theorem}

\begin{proof}
We only derive this form from Chow--Luo's results.  Let
\[
  K:\R_{>0}^V\longrightarrow \R^V,\qquad r\longmapsto K(r)
\]
be the hyperbolic-background curvature map.  Chow--Luo's description of
the curvature image says that a vector $x\in\R^V$ lies in $K(\R_{>0}^V)$
if and only if, for every nonempty proper subset $I\subsetneq V$,
\[
  \sum_{i\in I}x_i
  >
  -\sum_{(e,v)\in\Lk(I)}(\pi-\Phiw(e))+2\pi\chi(F_I).
\]
Applying this to $x=0$ gives exactly
\[
  \sum_{(e,v)\in\Lk(I)}(\pi-\Phiw(e))-2\pi\chi(F_I)>0,
\]
that is, $B_\Phi(I)>0$, for every nonempty proper $I$.

The additional full-surface condition is Gauss--Bonnet.  In the hyperbolic
background,
\[
  \sum_{i\in V}K_i(r)=2\pi\chi(X)+\operatorname{Area}(X,r).
\]
Thus if $K(r)=0$, then
\[
  \operatorname{Area}(X,r)=-2\pi\chi(X)>0,
\]
so $\chi(X)<0$.  Conversely, if $\chi(X)<0$ and all the proper-subset
inequalities $B_\Phi(I)>0$ hold, then the curvature image theorem applied
to $x=0$ gives a radius vector $r_{ze}\in\R_{>0}^V$ with $K(r_{ze})=0$.

Finally, Chow--Luo's convergence theorem for the hyperbolic combinatorial
Ricci flow states that, when such a zero-curvature hyperbolic circle
packing exists, the flow
\[
  \frac{dr_i}{dt}=-K_i\sinh r_i
\]
exists for all time and converges exponentially fast to it from every
initial radius vector in $\R_{>0}^V$.  Uniqueness follows from the
rigidity/injectivity theorem for the hyperbolic-background curvature map.
\end{proof}

\begin{theorem}[average character criterion]\label{thm:average-character-criterion}
\uses{thm:chow-luo,lem:obstruction-corrected-surplus,def:boundary-corrected-average}
\lean{ProbabilityCharacter.average_character_criterion_skeleton, ProbabilityCharacter.obstruction_pos_of_correctedAverage_gt, ProbabilityCharacter.eulerChar_neg_of_average_formula}
\leanok
Let $(X,\T,\Phiw)$ be a weighted triangulated closed surface with
$\Phiw:E\to[0,\pi/2]$.  Suppose that
\[
  \ac_\Phi(I)>2\pi
  \qquad
  \text{for every nonempty } I\subseteq V .
\]
Then $\chi(X)<0$, and the hyperbolic combinatorial Ricci flow
\[
  \frac{dr_i}{dt}=-K_i\sinh r_i
\]
exists for all time and converges exponentially fast, from every initial
radius vector $r(0)\in\R_{>0}^{V}$, to the unique zero-curvature
hyperbolic circle packing based on $(X,\T,\Phiw)$.
\end{theorem}

\begin{proof}
Taking $I=V$ gives
\[
  \frac{1}{|V|}\sum_{i\in V}L_i>2\pi.
\]
By Ge--Lin's average character formula,
\[
  \frac{1}{|V|}\sum_{i\in V}L_i
  =
  2\pi\left(1-\frac{\chi(X)}{|V|}\right),
\]
and therefore $\chi(X)<0$.

For each nonempty proper $I\subsetneq V$, the obstruction identity gives
\[
  B_\Phi(I)=|I|\bigl(\ac_\Phi(I)-2\pi\bigr)>0.
\]
Thus the zero curvature vector satisfies all Chow--Luo image inequalities.
The conclusion follows from Theorem~\ref{thm:chow-luo}.
\end{proof}

\begin{corollary}[Ge--Lin's pointwise criterion]\label{cor:ge-lin-pointwise}
\uses{thm:average-character-criterion,lem:nonneg-correction}
\lean{ProbabilityCharacter.pointwise_gt_two_pi_correctedAverage_gt, ProbabilityCharacter.pointwise_criterion_implies_average_criterion}
\leanok
If
\[
  L_i(\T,\Phiw)>2\pi
  \qquad (i\in V),
\]
then the conclusion of the average character criterion holds.
\end{corollary}

\begin{proof}
For $I=V$, the pointwise hypothesis gives $\ac_\Phi(V)>2\pi$.  For every
nonempty proper $I\subsetneq V$,
\[
  L_\Phi(I)=\sum_{i\in I}L_i>2\pi |I|,
\]
and $\Delta_\Phi(I)\ge0$ by the non-negativity lemma.  Hence
$\ac_\Phi(I)>2\pi$ for all nonempty $I\subseteq V$, and the theorem
applies.
\end{proof}

\begin{corollary}[lower average character criterion]\label{cor:lower-average-character-criterion}
\uses{thm:average-character-criterion,lem:obstruction-corrected-surplus,lem:nonneg-correction}
\lean{ProbabilityCharacter.lower_average_character_criterion}\leanok
Assume $\chi(X)<0$.  Suppose that there are constants $\eta>0$ and
$M\ge1$ such that
\[
  \frac{1}{|I|}\sum_{i\in I}L_i\ge2\pi+\eta
\]
for every proper connected subset $I\subsetneq V$ with $|I|\ge M$, and
that
\[
  \ac_\Phi(I)>2\pi
\]
for every proper connected subset $I\subsetneq V$ with $|I|<M$.  Then the
hyperbolic combinatorial Ricci flow converges exponentially fast to the
unique zero-curvature hyperbolic circle packing.
\end{corollary}

\begin{proof}
It is enough to verify $\ac_\Phi(I)>2\pi$ for every nonempty proper
$I\subsetneq V$.  Let such an $I$ be given, and decompose the full
subcomplex $F_I$ into connected components $F_{I_1},\dots,F_{I_s}$, where
$I=I_1\sqcup\cdots\sqcup I_s$ and each $I_a$ is connected in the induced
one-skeleton.  A triangle with at least two vertices in $I$ has those two
vertices joined by an edge, hence they lie in the same component.  Therefore
the terms defining $B_\Phi(I)$ split over the components:
\[
  B_\Phi(I)=\sum_{a=1}^s B_\Phi(I_a).
\]

If $|I_a|\ge M$, then
\[
  L_\Phi(I_a)>2\pi |I_a|
\]
and $\Delta_\Phi(I_a)\ge0$, so $B_\Phi(I_a)>0$ by the obstruction identity.
If $|I_a|<M$, then the assumed corrected inequality gives
$B_\Phi(I_a)>0$.  Thus $B_\Phi(I)>0$ for every nonempty proper $I$.
The assumption $\chi(X)<0$ gives $\ac_\Phi(V)>2\pi$ by Ge--Lin's average
character formula.  The average character criterion applies.
\end{proof}

\section{Usable forms of the criterion}

The quantity $\ac_\Phi(I)$ is useful for the proof because it is exactly
the Chow--Luo obstruction written in character language.  It is not the
most convenient form for applications.  In this section we extract simpler
tests from it.

\begin{corollary}[only potentially bad subsets need checking]\label{cor:only-potentially-bad-subsets}
\uses{cor:lower-average-character-criterion,thm:average-character-criterion,lem:nonneg-correction}
\lean{ProbabilityCharacter.corollary_one_only_potentially_bad_subsets}\leanok
Assume $\chi(X)<0$.  Suppose that
\[
  \ac_\Phi(I)>2\pi
\]
for every nonempty proper connected subset $I\subsetneq V$ satisfying
\[
  L_\Phi(I)\le 2\pi |I|.
\]
Then the hyperbolic combinatorial Ricci flow converges exponentially fast
to the unique zero-curvature hyperbolic circle packing.
\end{corollary}

\begin{proof}
Let $I\subsetneq V$ be nonempty and connected.  If
$L_\Phi(I)>2\pi |I|$, then $\Delta_\Phi(I)\ge0$ gives
$\ac_\Phi(I)>2\pi$.  If $L_\Phi(I)\le2\pi |I|$, the desired inequality is
assumed.  Hence the corrected average inequality holds for every nonempty
proper connected subset.  The proof of the lower average character
criterion shows that this is enough to get it for arbitrary nonempty proper
subsets, and then the average character criterion applies.
\end{proof}

\begin{corollary}[constant weight formula]\label{cor:constant-weight-formula}
\uses{def:boundary-corrected-average,lem:comparison-angle}
\lean{ProbabilityCharacter.constant_weight_formula, ProbabilityCharacter.constant_weight_formula_nat}\leanok
Assume $\Phiw\equiv\varphi$ is constant.  For every nonempty proper
subset $I\subsetneq V$,
\[
  \ac_\Phi(I)>2\pi
\]
is equivalent to
\[
  \sum_{i\in I}(d_i-6)+|\mathcal F_2(I)|
  +\left(2-\frac{3\varphi}{\pi}\right)|\mathcal F_1(I)|>0.
\]
\end{corollary}

\begin{proof}
For constant weight, each comparison angle is $\pi/3$.  Hence
\[
  L_\Phi(I)=\frac{\pi}{3}\sum_{i\in I}d_i
\]
and
\[
  \Delta_\Phi(I)
  =
  \frac{\pi}{3}|\mathcal F_2(I)|
  +\left(\frac{2\pi}{3}-\varphi\right)|\mathcal F_1(I)|.
\]
The inequality $L_\Phi(I)+\Delta_\Phi(I)>2\pi |I|$ is therefore equivalent
to the displayed formula after multiplying by $3/\pi$.
\end{proof}

\begin{corollary}[tangent circle packings]\label{cor:tangent-circle-packings}
\uses{cor:only-potentially-bad-subsets,cor:constant-weight-formula}
\lean{ProbabilityCharacter.corollary_two_tangent_circle_packings}\leanok
Assume $\Phiw\equiv0$ and $\chi(X)<0$.  Suppose that for every nonempty
proper connected subset $I\subsetneq V$ with
\[
  \sum_{i\in I}d_i\le 6|I|
\]
one has
\[
  \sum_{i\in I}(d_i-6)+|\mathcal F_2(I)|+2|\mathcal F_1(I)|>0.
\]
Then the hyperbolic combinatorial Ricci flow converges exponentially fast
to the unique zero-curvature tangent circle packing.
\end{corollary}

\begin{proof}
This is the previous two corollaries in the special case $\varphi=0$.
\end{proof}

\begin{remark}[bad clusters and discharging]
The preceding tangent criterion is often much easier to check than the
exact $\ac_\Phi$ inequality: only connected subsets whose average degree is
at most six can be dangerous, and their boundary faces contribute the
positive correction
\[
  |\mathcal F_2(I)|+2|\mathcal F_1(I)|.
\]
This suggests a local ``bad cluster'' strategy.  Let the initial charge be
$d_i-6$ in the tangent case, or more generally $L_i-2\pi$.  Low-charge
vertices may be allowed if every connected cluster with non-positive total
charge receives enough boundary correction.  A still more local
discharging criterion would distribute the positive surplus and boundary
correction among nearby low-charge vertices.  Such a discharging lemma
would be a convenient sufficient condition for the exact criterion above,
but it should be stated separately once the chosen redistribution rule is
fixed.
\end{remark}

\section{Degree criteria as special cases}

The corrected average character criterion contains the familiar
degree-type criteria as special cases.

\begin{corollary}[constant weight, pointwise degree criterion]\label{cor:constant-weight-pointwise-degree}
\uses{cor:ge-lin-pointwise,cor:constant-weight-formula}
\lean{ProbabilityCharacter.pointwise_degree_criterion_constant_weight}\leanok
Assume that $\Phiw:E\to[0,\pi/2]$ is constant.  If
\[
  d_i\ge7\qquad (i\in V),
\]
then the hyperbolic combinatorial Ricci flow converges exponentially fast
to the unique zero-curvature hyperbolic circle packing.
\end{corollary}

\begin{proof}
When $\Phiw$ is constant, Ge--Lin's character at a vertex is
\[
  L_i=\sum_{f\ni i}\frac{\pi}{3}=\frac{\pi d_i}{3}.
\]
Thus $d_i\ge7$ implies
\[
  L_i\ge\frac{7\pi}{3}>2\pi
\]
for every vertex.  The pointwise character corollary applies.
\end{proof}

\begin{corollary}[arbitrary weight, pointwise degree criterion]\label{cor:arbitrary-weight-pointwise-degree}
\uses{cor:ge-lin-pointwise}
\lean{ProbabilityCharacter.pointwise_degree_criterion_arbitrary_weight}\leanok
Assume that $\Phiw:E\to[0,\pi/2]$ is arbitrary.  If
\[
  d_i\ge9\qquad (i\in V),
\]
then the hyperbolic combinatorial Ricci flow converges exponentially fast
to the unique zero-curvature hyperbolic circle packing.
\end{corollary}

\begin{proof}
For $\Phiw\in[0,\pi/2]$, Ge--Lin's angle estimate gives
\[
  L_i\ge d_i\arccos\frac34.
\]
Since $9\arccos(3/4)>2\pi$, the condition $d_i\ge9$ implies
$L_i>2\pi$ for every vertex.  The pointwise character corollary applies.
\end{proof}

\begin{corollary}[constant weight, large-set lower average degree criterion]\label{cor:large-set-lower-average-degree}
\uses{cor:lower-average-character-criterion,cor:constant-weight-formula}
\lean{ProbabilityCharacter.corollary_three_large_set_lower_average_degree}\leanok
Assume that $\Phiw:E\to[0,\pi/2]$ is constant and that $\chi(X)<0$.
Suppose that there are a constant $\eta>0$ and an integer $M\ge2$ such that
\[
  \frac1{|I|}\sum_{i\in I}d_i \ge 6+\eta
\]
for every proper connected subset $I\subsetneq V$ with $|I|\ge M$, and
that every smaller proper connected subset satisfies the corrected test
\[
  \ac_\Phi(I)>2\pi
\]
whenever $|I|<M$.  Then the hyperbolic combinatorial Ricci flow converges
exponentially fast to the unique zero-curvature hyperbolic circle packing.
\end{corollary}

\begin{proof}
For constant weight, $L_i=\pi d_i/3$.  Hence, for every proper connected
$I\subsetneq V$ with $|I|\ge M$,
\[
  \frac1{|I|}\sum_{i\in I}L_i
  =
  \frac{\pi}{3|I|}\sum_{i\in I}d_i
  \ge
  2\pi+\frac{\pi\eta}{3}.
\]
The lower average character criterion applies, with $\pi\eta/3$ in place
of $\eta$.
\end{proof}

\begin{remark}[the singleton test does not force degree at least seven]
The restriction $M\ge2$ is not needed for the proof; the statement would
remain true with $M=1$, but then the large-set average degree condition
would apply to singletons and would force $d_v\ge7$ in the constant
tangency case.  The present formulation is the intended large-set version:
singletons are handled by the corrected small-set test instead.

For a singleton $I=\{v\}$ and constant weight $\Phiw\equiv\varphi$, every
incident face lies in $\mathcal F_1(I)$ and
\[
  L_v=\frac{\pi d_v}{3},\qquad
  \Delta_\Phi(\{v\})
  =
  d_v\left(\pi-\varphi-\frac{\pi}{3}\right).
\]
Therefore
\[
  \ac_\Phi(\{v\})=d_v(\pi-\varphi).
\]
In particular, in the tangency case $\varphi=0$, the singleton condition is
only $d_v\pi>2\pi$, i.e. $d_v>2$.  Thus vertices of degree $3,4,5,$ or $6$
are not excluded by the singleton condition.  Low-degree vertices are
allowed provided the corrected inequalities for all small connected
subsets containing them are satisfied.
\end{remark}

\section{Relation to the source ideas}

The result above is a finite closed-surface analogue of the lower average
valence idea in He--Schramm: a pointwise local condition is replaced by
conditions over finite connected vertex sets.  The quantity replacing
average degree is the corrected average character $\ac_\Phi(I)$.

It also mirrors the angle-sum bookkeeping in
Angel--Hutchcroft--Nachmias--Ray.  There, mass transport converts local
angle sums of circle-packed triangles into an average degree criterion.  In
the finite setting here, the same aggregation appears as the exact identity
\[
  B_\Phi(I)=|I|(\ac_\Phi(I)-2\pi),
\]
which is the boundary-corrected finite subset version of the same angle
defect bookkeeping.
